\section{Unconditional cancellation for regular plane kernels}
\label{sec:smooth-kernels}

The master certificate has a positive unconditional consequence.  It is
useful both as a theorem and as a calibration of what the physical kernel
would need to resemble.

\begin{theorem}[Bounded fixed-d variation implies cancellation]
\label{thm:regular-kernel-cancellation}
Let \(K_{U,V}\) be supported on
\([1,U]\times[1,V]\), and put
\[
 B_{U,V}=\|K_{U,V}\|_\infty.
\]
Assume \(B_{U,V}>0\) and, for every fixed \(d\ge1\),
\begin{equation}
 \cV_\square(K_{U,V}^{[d]})
 \le C_d B_{U,V}
\label{eq:fixed-d-variation-hypothesis}
\end{equation}
with \(C_d\) independent of \(U,V\).  Then, as
\(\min(U,V)\to\infty\),
\begin{equation}
 \boxed{
 \sum_{\substack{u\le U,\ v\le V\\(u,v)=1}}
 \mu(u)\mu(v)K_{U,V}(u,v)
 =
 o(B_{U,V}UV).}
\label{eq:regular-kernel-cancellation}
\end{equation}
\end{theorem}

\begin{proof}
Fix \(D\).  By \cref{lem:fixed-d-mertens}, for every \(d\le D\),
\[
 \mathfrak M_d(U/d)\mathfrak M_d(V/d)
 =
 o_d(UV/d^2).
\]
The finitely many low-diagonal terms in
\eqref{eq:master-certificate}, divided by \(B_{U,V}UV\), therefore tend to
zero by \eqref{eq:fixed-d-variation-hypothesis}.  The ambient tail
\eqref{eq:ambient-tail} contributes at most \(1/D\).  Take the scale limit
first and then \(D\to\infty\).
\end{proof}

\begin{corollary}[The primitive rectangle]
\label{cor:primitive-rectangle}
As \(\min(U,V)\to\infty\),
\begin{equation}
 \sum_{\substack{u\le U,\ v\le V\\(u,v)=1}}
 \mu(u)\mu(v)=o(UV).
\label{eq:primitive-rectangle}
\end{equation}
\end{corollary}

\begin{proof}
Take \(K=\one_{[1,U]\times[1,V]}\).  For every \(d\), its zero-extended
dilation is another rectangle and
\[
 \cV_\square(K^{[d]})=1.
\]
Apply \cref{thm:regular-kernel-cancellation}.
\end{proof}

\begin{remark}
\Cref{cor:primitive-rectangle} is not the two-point Chowla conjecture.
The two variables are independently averaged over a full rectangle, and
the coprimality condition turns the sign into \(\mu(uv)\).  No prescribed
affine relation such as \(v=u+1\) is imposed.
\end{remark}

\subsection{Projective bounded variation}

For a finitely supported one-dimensional sequence \(f\), define its
zero-extended variation
\[
 \cV_0(f)=\sum_{x\ge1}|f(x)-f(x+1)|.
\]

\begin{lemma}[Projective variation bound]
\label{lem:projective-variation}
If
\[
 K(x,y)=\sum_{r=1}^R f_r(x)g_r(y),
\]
then
\begin{equation}
 \cV_\square(K)
 \le
 \sum_{r=1}^R\cV_0(f_r)\cV_0(g_r).
\label{eq:projective-variation}
\end{equation}
\end{lemma}

\begin{proof}
For a rank-one kernel,
\[
 \Delta_{12}(fg)(x,y)
 =
 \bigl(f(x)-f(x+1)\bigr)
 \bigl(g(y)-g(y+1)\bigr).
\]
Sum absolute values and then use the triangle inequality over \(r\).
\end{proof}

\begin{corollary}[Finite projective-BV families]
\label{cor:projective-bv}
Suppose that, for every fixed \(d\), the dilated kernels admit
representations
\[
 K_{U,V}^{[d]}(x,y)
 =
 \sum_r f_{r,d}(x)g_{r,d}(y)
\]
such that
\[
 \sum_r\cV_0(f_{r,d})\cV_0(g_{r,d})
 \ll_d B_{U,V}.
\]
Then \eqref{eq:regular-kernel-cancellation} holds.
\end{corollary}

\begin{proof}
Combine \cref{lem:projective-variation} with
\cref{thm:regular-kernel-cancellation}.
\end{proof}

\subsection{Dyadic localization}

The same argument applies to translated rectangles.  If \(K\) is supported
on
\[
 U_0<u\le U_1,\qquad V_0<v\le V_1,
\]
zero extension automatically inserts the four boundary jumps.  One may
also use interval partial sums
\[
 M_d(x)-M_d(U_0/d)
\]
inside the Abel formula.  On dyadic boxes with \(U,V\to\infty\),
\[
 U<u\le2U,\qquad V<v\le2V,
\]
fixed-\(d\) Mertens cancellation is still \(o(U)\) and \(o(V)\).
Therefore \cref{thm:regular-kernel-cancellation} is stable under a bounded
number of dyadic localizations.  Summing \(O(\log^2 X)\) boxes is harmless
only when the target ledger permits that subpower bookkeeping and when the
variation bounds are uniform over the boxes.
Boxes with \(U=O(1)\) or \(V=O(1)\) belong instead to the boundary or
anisotropic sectors of \cref{sec:physical-gate}.

\subsection{A squarefree-pair calibration}

The set on which \(\mu(uv)\ne0\) has positive planar density
\begin{equation}
 \mathfrak S_{\rm sf}
 =
 \prod_p\left(1-\frac3{p^2}+\frac2{p^3}\right)>0.
\label{eq:squarefree-pair-density}
\end{equation}
Indeed, the local condition is that the sum of the two \(p\)-adic
valuations is at most one.  Its local density is
\[
 (1-p^{-1})^2
 +2(p^{-1}-p^{-2})(1-p^{-1})
 =
 1-\frac3{p^2}+\frac2{p^3}.
\]
For a fixed \(z\), the Chinese remainder theorem therefore gives the
product of these factors over \(p\le z\), up to
\(O_z(U+V+1)\) boundary pairs.  A violation caused by a prime \(p>z\)
requires
\[
 p^2\mid u,\qquad\text{or}\qquad
 p^2\mid v,\qquad\text{or}\qquad p\mid(u,v).
\]
The union of these events contains at most
\[
 3UV\sum_{p>z}\frac1{p^2}
\]
pairs.  Let first \(\min(U,V)\to\infty\) and then \(z\to\infty\).
This proves
\begin{equation}
 \sum_{u\le U,v\le V}\mu(uv)^2
 =
 \mathfrak S_{\rm sf}UV+o(UV)
\label{eq:squarefree-pair-asymptotic}
\end{equation}
as \(\min(U,V)\to\infty\).  Together with
\cref{cor:primitive-rectangle}, this shows that the signs \(+1\) and
\(-1\) of \(\mu(uv)\) each occupy asymptotically
\(\mathfrak S_{\rm sf}UV/2\) points.  This fact will calibrate the hard-mask
countermodel in \cref{sec:countermodels}.
